% chapters/bab3-metodologi-penelitian.tex -- BAB 3 METODOLOGI PENELITIAN
% Contoh teks ilmiah di bawah ini adalah placeholder untuk memperlihatkan
% hasil format (thesis.md); ganti dengan isi penelitian yang sebenarnya.

\chapter{METODOLOGI PENELITIAN}

\section{Jenis dan Pendekatan Penelitian}
Penelitian ini menggunakan pendekatan eksperimental kuantitatif, yaitu
dengan membangun model klasifikasi sentimen berbasis LSTM,
melatihnya menggunakan data ulasan produk berbahasa Indonesia, dan
mengukur kinerjanya menggunakan metrik evaluasi standar pada bidang
pembelajaran mesin. Hasil model yang diusulkan kemudian dibandingkan
dengan model \textit{baseline} berbasis Support Vector Machine (SVM)
untuk mengetahui sejauh mana peningkatan kinerja yang dapat dicapai.

\section{Data Penelitian}
Data yang digunakan dalam penelitian ini adalah kumpulan ulasan
produk berbahasa Indonesia yang diambil dari salah satu platform
perdagangan elektronik, terdiri atas teks ulasan beserta label
sentimen (positif, negatif, netral) yang diberikan oleh pengguna
melalui penilaian bintang. Data tersebut dibagi menjadi data latih,
data validasi, dan data uji dengan proporsi 70:15:15 untuk memastikan
evaluasi model dilakukan pada data yang belum pernah dilihat selama
proses pelatihan.

\section{Tahapan Penelitian}
Penelitian ini dilaksanakan melalui beberapa tahapan sebagai berikut.
\begin{enumerate}
  \item \textbf{Pengumpulan data.} Data ulasan produk dikumpulkan
        beserta label sentimennya, kemudian diperiksa kualitasnya
        untuk memastikan tidak ada data yang duplikat atau tidak
        relevan.
  \item \textbf{Praproses data.} Data ulasan melalui tahapan
        \textit{case folding}, tokenisasi, normalisasi kata tidak
        baku, penghapusan stopword, dan \textit{stemming}
        menggunakan kamus normalisasi bahasa Indonesia.
  \item \textbf{Ekstraksi fitur.} Teks hasil praproses diubah menjadi
        representasi vektor menggunakan teknik \textit{word
        embedding}, yang kemudian digunakan sebagai masukan bagi
        lapisan LSTM.
  \item \textbf{Perancangan dan pelatihan model.} Model LSTM
        dirancang dan dilatih menggunakan data latih, dengan
        penyesuaian hiperparameter dilakukan berdasarkan kinerja pada
        data validasi.
  \item \textbf{Evaluasi model.} Kinerja model diuji menggunakan data
        uji dan dibandingkan dengan model \textit{baseline}
        menggunakan metrik evaluasi yang telah ditentukan.
\end{enumerate}

\section{Perancangan Model}

\subsection{Arsitektur model LSTM yang diusulkan}
Model yang diusulkan pada penelitian ini terdiri atas lapisan
\textit{embedding}, satu lapisan LSTM dengan 128 unit tersembunyi,
lapisan \textit{dropout} untuk mengurangi risiko \textit{overfitting},
dan lapisan \textit{fully connected} dengan fungsi aktivasi softmax
pada lapisan keluaran untuk menghasilkan probabilitas masing-masing
kelas sentimen. Ringkasan arsitektur model disajikan pada
Tabel~\ref{tab:arsitektur}.

\begin{table}[htbp]
  \centering
  \caption{Ringkasan arsitektur model LSTM yang diusulkan}
  \label{tab:arsitektur}
  \singlespaced{%
    \begin{tabular}{lll}
      \hline
      Lapisan & Jumlah Unit & Fungsi Aktivasi \\
      \hline
      Embedding       & 128 & -- \\
      LSTM             & 128 & tanh \\
      Dropout (0.5)    & --  & -- \\
      Fully Connected  & 3   & softmax \\
      \hline
    \end{tabular}%
  }
\end{table}

\subsection{Penyesuaian hiperparameter}
Penyesuaian hiperparameter dilakukan terhadap ukuran \textit{batch},
laju pembelajaran (\textit{learning rate}), dan jumlah epoch
pelatihan. Nilai hiperparameter terbaik dipilih berdasarkan kinerja
tertinggi pada data validasi, guna menghindari model yang terlalu
menyesuaikan diri pada data latih.

\section{Metode Evaluasi}
Kinerja model dievaluasi menggunakan metrik akurasi, presisi, recall,
dan F1-score, yang dihitung berdasarkan matriks konfusi hasil
klasifikasi pada data uji. Perhitungan F1-score dapat dituliskan
sebagai berikut.
\begin{equation}
  \text{F1} = 2 \times \frac{\text{presisi} \times \text{recall}}{\text{presisi} + \text{recall}}
  \label{eq:f1score}
\end{equation}
Persamaan~\eqref{eq:f1score} digunakan karena F1-score memberikan
gambaran yang lebih seimbang dibandingkan akurasi semata, terutama
apabila distribusi kelas pada data uji tidak seimbang.
